%*************************************
% บทที่ 5
%*************************************
\chapterTitle{5}{สรุปผลการดำเนินงาน}

\section{สรุปผลการดำเนินงาน}
\hspace*{1.3em}
โครงการ "Cloud-Based Platform for Supporting Teaching and Academic Activities" ได้พัฒนาแพลตฟอร์มคลาวด์สำหรับสนับสนุนการเรียนการสอนและกิจกรรมทางวิชาการในภาควิชาเทคโนโลยีสารสนเทศให้สามารถใช้งานได้จริงในลักษณะ Microservices Architecture โดยมีบริการหลัก 3 ส่วน ได้แก่ Midori สำหรับส่วนติดต่อผู้ใช้ Momoi สำหรับ API และตรรกะทางธุรกิจ และ Yuzu สำหรับการประมวลผลงานจัดการเครื่องเสมือน ร่วมกับบริการสนับสนุนด้านการยืนยันตัวตน ฐานข้อมูล ระบบคิวงาน พื้นที่จัดเก็บไฟล์ และระบบติดตามการทำงาน

\hspace*{1.3em}
ผลลัพธ์สำคัญของโครงการคือการสร้างแพลตฟอร์มที่รองรับวงจรการใช้งานเครื่องเสมือนตั้งแต่การยื่นคำร้อง การอนุมัติ การจัดสรรทรัพยากร การเปิดใช้งาน การควบคุมสถานะ ไปจนถึงการลบคืนทรัพยากร โดยเชื่อมโยงกระบวนการดังกล่าวเข้ากับข้อมูลทางวิชาการ เช่น รายวิชา ภาคการศึกษา และผู้สอน ทำให้การให้บริการเครื่องเสมือนสอดคล้องกับบริบทการเรียนการสอนจริงของภาควิชา

\hspace*{1.3em}
นอกจากนี้ ระบบยังรองรับการควบคุมสิทธิ์แบบ Role-Based Access Control (RBAC) สำหรับผู้ใช้กลุ่ม STUDENT, INSTRUCTOR และ ADMIN มีเอกสาร API แบบ OpenAPI สำหรับการพัฒนาและทดสอบ มีส่วนจัดการไฟล์บน S3-compatible storage และมีชุด observability สำหรับติดตาม metrics, logs และ traces ของบริการหลัก ผลการดำเนินงานโดยรวมจึงสะท้อนว่าระบบสามารถเป็นพื้นฐานของแพลตฟอร์มดิจิทัลสำหรับสนับสนุนการเรียนการสอนและกิจกรรมทางวิชาการได้อย่างเป็นรูปธรรม

\hspace*{1.3em}
ในด้านสมรรถนะ ผลการทดสอบประสิทธิภาพบริการหลักพบว่า Midori รองรับคำขอได้ประมาณ 720 คำขอต่อวินาที (RPS) ด้วยอัตราความสำเร็จ 100\% และ latency median ราว 54–56 ms สำหรับหน้าหลัก ขณะที่ Momoi รองรับได้ประมาณ 2{,}255 RPS ด้วย latency median ราว 16 ms สำหรับ API สาธารณะ ส่วนการทดสอบ Yuzu พบว่า Linked Clone มีอัตราความสำเร็จ 100\% และเร็วกว่า Full Clone โดยเฉลี่ยกว่า 96\% ผลการทดสอบดังกล่าวบ่งชี้ว่าระบบมีประสิทธิภาพเพียงพอสำหรับการใช้งานจริงในบริบทการเรียนการสอนของภาควิชา

\section{ปัญหาและอุปสรรค}
\begin{mycustomenum}[label=5.2.\arabic*] 
    \item ความซับซ้อนของการประสานงานระหว่างบริการหลายส่วน เนื่องจากระบบประกอบด้วย frontend, API, worker และบริการสนับสนุนจำนวนมากที่ต้องทำงานร่วมกันผ่าน HTTP API, Redis และฐานข้อมูล ทำให้การดีบักและติดตามปัญหาต้องอาศัยความเข้าใจภาพรวมของระบบ
    \item การจัดการเครื่องเสมือนผ่าน Proxmox VE API มีรายละเอียดเชิงเทคนิคสูง ทั้งด้านสถานะของ VM, การเลือกโหนด, การจัดสรรเครือข่าย และการจัดคิวงานที่ใช้เวลานาน โดยเฉพาะงานประเภท clone หรือ reprovision ที่มีผลต่อทรัพยากรของระบบโดยตรง
    \item การออกแบบโมเดลข้อมูลให้รองรับทั้งบริบททางวิชาการและการควบคุมสิทธิ์ของผู้ใช้เป็นโจทย์สำคัญ เนื่องจากระบบไม่ได้จัดการเพียงบัญชีผู้ใช้ แต่ยังต้องเชื่อมโยงข้อมูลรายวิชา ภาคการศึกษา ผู้สอน คำร้อง และเครื่องเสมือนเข้าด้วยกันอย่างถูกต้อง
    \item การใช้เทคโนโลยีใหม่ในหลายชั้นของระบบ เช่น Bun Runtime, Elysia.js, BullMQ, Prisma และแนวทางการสร้าง type-safe client จาก OpenAPI ต้องใช้เวลาในการศึกษา ทดลอง และกำหนดแนวทางที่เหมาะสมกับโครงงาน
    \item การติดตามประสิทธิภาพและความผิดปกติของระบบในสภาพแวดล้อมที่มีหลายบริการยังเป็นความท้าทาย โดยเฉพาะเมื่อต้องวิเคราะห์ความล่าช้าจากหลายจุด เช่น frontend proxy, API handler, queue worker และการเรียกใช้งาน Proxmox VE
\end{mycustomenum}

\section{การแก้ปัญหา}
\begin{mycustomenum}[label=5.3.\arabic*] 
    \item แยกความรับผิดชอบของระบบออกเป็นรีโพซิทอรีและบริการที่ชัดเจน ได้แก่ Midori, Momoi, Yuzu และ Yuuka พร้อมใช้ Dockerfile และไฟล์ deployment เพื่อให้สามารถพัฒนา ทดสอบ และติดตั้งใช้งานแต่ละส่วนได้อย่างเป็นอิสระ
    \item ใช้ OpenAPI เป็นมาตรฐานสำหรับการอธิบาย API ของ Momoi และนำข้อมูลดังกล่าวไปใช้สร้างชนิดข้อมูลฝั่ง client ช่วยลดความคลาดเคลื่อนระหว่าง frontend และ backend รวมถึงช่วยให้เอกสารประกอบการพัฒนามีความชัดเจนมากขึ้น
    \item พัฒนา Yuzu ให้ทำงานเป็น background worker ที่รับงานจาก Redis-backed BullMQ แทนการประมวลผลแบบ synchronous ภายใน API โดย Momoi จะส่งเพียงข้อมูลอ้างอิงที่จำเป็นเข้าสู่คิว จากนั้น Yuzu จึงอ่านข้อมูลเพิ่มเติมจากฐานข้อมูลก่อนดำเนินการกับ Proxmox VE ทำให้การ retry และการจัดการงานระยะยาวทำได้เหมาะสมขึ้น
    \item ออกแบบฐานข้อมูลด้วย Prisma ORM ให้รองรับทั้งข้อมูลผู้ใช้ โครงสร้างวิชาการ คำร้องขอ เครื่องเสมือน reverse proxy และพื้นที่จัดเก็บไฟล์ พร้อมกำหนดบทบาทของผู้ใช้อย่างชัดเจน และเชื่อมโยงการยืนยันตัวตนกับบริการ auth ที่ใช้งานร่วมกับแพลตฟอร์ม
    \item ใช้ OpenTelemetry และติดตั้งชุด observability เช่น Prometheus, Grafana, Loki และ Jaeger เพื่อเก็บ metrics, logs และ traces ของแต่ละบริการ ช่วยให้สามารถวิเคราะห์ปัญหาเชิงระบบและติดตามประสิทธิภาพได้มีประสิทธิภาพมากขึ้น
\end{mycustomenum}

\section{ข้อเสนอแนะ}
\begin{mycustomenum}[label=5.4.\arabic*] 
    \item พัฒนาระบบ auto-scaling, load balancing และการแยกทรัพยากรตามลักษณะงาน เพื่อรองรับจำนวนผู้ใช้งานที่เพิ่มขึ้น โดยเฉพาะในช่วงที่มีการสร้างเครื่องเสมือนพร้อมกันจำนวนมากจากกิจกรรมการเรียนการสอน
    \item เพิ่มความสามารถด้าน VM template lifecycle เช่น การจัดการ template ตามรายวิชา การทำ snapshot, backup และการกู้คืนสถานะ เพื่อให้รองรับการใช้งานจริงในรายวิชาที่มีความต้องการแตกต่างกันมากขึ้น
    \item ปรับปรุงกระบวนการ queue worker และการจัดการงานบนสตอเรจหรือคลัสเตอร์ให้รองรับการทำงานพร้อมกันได้ดีขึ้น โดยเฉพาะกรณี full clone หรือการดำเนินงานที่ใช้ทรัพยากรสูง ซึ่งจากผลการทดสอบพบว่ายังมีข้อจำกัดด้านเวลาและอัตราความสำเร็จเมื่อเทียบกับ linked clone
    \item เพิ่มระบบ CI/CD และขยายชุดการทดสอบอัตโนมัติให้ครอบคลุมมากขึ้น ทั้ง unit test, integration test, end-to-end test และ performance regression test เพื่อให้การพัฒนาและปรับปรุงระบบในอนาคตมีความมั่นคงมากขึ้น
    \item ต่อยอดการติดตั้งใช้งานบนแพลตฟอร์ม orchestration ระดับ production เช่น Kubernetes อย่างเต็มรูปแบบ พร้อมกำหนดนโยบายด้านความปลอดภัย การจัดการ secret การสำรองข้อมูล และ disaster recovery เพื่อเพิ่มความพร้อมสำหรับการใช้งานในระดับองค์กร
\end{mycustomenum}